\label{sec:measure-theory}

At their core, (large) language models are an attempt to place a probabilistic distribution over natural language utterances.
However, our toy examples in \cref{ex:inf-coin-toss,ex:possibly-inf-coin-toss} in the previous section reveal that it is quite nuanced to get a satisfying definition of a language model.
Thus, our first step forward is to review the basics of rigorous probability theory,\footnote{By rigorous probability theory we mean a measure-theoretic treatment of probability theory.} the tools we need to come to a satisfying definition. 
Our course will assume that you have had some exposure to rigorous probability theory before, and just review the basics.
However, it is also possible to learn the basics of rigorous probability on the fly during the course if it is new to you.
Specifically, we will cover \emph{measure-theoretic} foundations of probability theory.
This might come as a bit of a surprise since we are mostly going to be talking about discrete objects---namely, strings.
However, as we will see in \cref{sec:lm-tightness} soon, formal treatment of language modeling indeed requires some mathematical rigor from measure theory.

The goal of measure-theoretic probability is to assign probabilities to \emph{subsets}
of an \defn{outcome space} $\samplespace$. 
However, in the course of the study of measure theory, it has become clear that for many common $\samplespace$, it is impossible to assign probabilities in a way that satisfies a set of reasonable desiderata.\footnote{Measure theory texts commonly discuss such desiderata and the dilemma that comes with it. See, e.g., Chapter 7 in \citet{Tao2016}, Chapter 3 in \citet{Royden1988} or Chapter 3 in \citet{Billingsley1986}. 
We also give an example later.} 
Consequently, the standard approach to probability theory resorts to only assigning probability to certain ``nice'' (but not all) subsets of $\samplespace$, which are referred to as \defn{events} or \defn{measurable subsets}, as in the theory of integration or functional analysis.
The set of measurable subsets is commonly denoted as $\eventspace$ (\cref{def:sigma-algebra}) and a probability measure $\probfunction:\eventspace\to[0,1]$ is the function that assigns a probability to each measurable subset. 
The triple $(\samplespace,\eventspace,\probfunction)$ is collectively known as a probability space (\cref{def:probability-measure}).
As it turns out, the following simple and reasonable requirements imposed on $\eventspace$ and $\probfunction$ are enough to rigorously discuss probability.

\begin{definition} \label{def:sigma-algebra}
Let $\powerset{\samplespace}$ be the power set of $\samplespace$.
Then $\eventspace \subseteq \powerset\samplespace$ is called a \defn{$\sigalgebra$} (or $\sigma$-field) over $\samplespace$ if the following conditions hold:
\begin{enumerate}[label=\arabic*)]
    \item $\samplespace\in\eventspace$,
    \item if $\sE\in\eventspace$, then $\setcomplement{\sE} \in \eventspace$,
    \item if $\sE_1,\sE_2,\dots$ is a finite or infinite sequence of sets in $\eventspace$, then $\bigcup_n \sE_n \in\eventspace$.
\end{enumerate}
If $\eventspace$ is a $\sigma$-algebra over $\samplespace$, we call the tuple $(\samplespace,\eventspace)$ a \defn{measurable space}.
\end{definition}

\begin{example}[$\sigalgebra$s]
Let $\samplespace$ be any set. 
Importantly, there is no one way to construct a $\sigalgebra$ over $\samplespace$:
\begin{enumerate}
    \item The family consisting of only the empty set $\emptyset$ and the set $\samplespace$, i.e., $\eventspace\defeq \left\{\emptyset, \samplespace\right\}$, is called the \emph{minimal} or \emph{trivial}.
    \item The full power set $\eventspace \defeq \powerset{\samplespace}$ is called the \emph{discrete $\sigalgebra$}.
    \item Given $\sA \subseteq \samplespace$, the family $\eventspace \defeq \left\{\emptyset, \sA, \samplespace \setminus \sA, \samplespace \right\}$ is a $\sigalgebra$ induced by $\sA$.
    \item Suppose we are rolling a six-sided die. 
    There are six events that can happen: we can roll any of the numbers $1\text{--}6$.
    In this case, we will then define the set of outcomes $\samplespace$ as $\samplespace \defeq \left\{\text{The number observed is } n \mid n = 1, \ldots, 6\right\}$.
    There are of course multiple ways to define an event space $\eventspace$ and with it a $\sigalgebra$ over this sample space.
    By definition, $\emptyset \in \eventspace$ and $\samplespace \in \eventspace$.
    One way to intuitively construct a $\sigalgebra$ is to consider that all individual events (observing any number) are possible, meaning that we would like to later assign probabilities to them (see \cref{def:probability-measure}). 
    This means that we should include individual singleton events in the event space: $\left\{\text{The number observed is } n \right\} \in \eventspace$ for $n =1, \ldots, 6$.
    It is easy to see that in this case, to satisfy the axioms in \cref{def:sigma-algebra}, the resulting event space should be $\eventspace = \powerset{\samplespace}$. 
\end{enumerate}
You might want to confirm these are indeed $\sigalgebra$s by checking them against the axioms in \cref{def:sigma-algebra}.
\end{example}
A measurable space guarantees that operations on countably many sets are always valid, and hence permits the following definition.
\begin{definition}
\label{def:probability-measure}
A \defn{probability measure} $\probfunction$ over a measurable space $(\samplespace,\eventspace)$ is a function $\probfunction:\eventspace\to[0,1]$ such that
\begin{enumerate}[label=\arabic*)]
    \item $\probfunction(\samplespace)=1$,
    \item if $\sE_1,\sE_2,\dots$ is a countable sequence of disjoint sets in $\eventspace$, then $\probfunction(\bigcup_n \sE_n)=\sum_n \probfunction(\sE_n)$.
\end{enumerate}
In this case we call $(\samplespace,\eventspace, \probfunction)$ a \defn{probability space}.\looseness=-1
\end{definition}

As mentioned, measure-theoretic probability only assigns probabilities to ``nice'' subsets of $\samplespace$.
In fact, it is often impossible to assign a probability measure to every single subset of $\samplespace$ and we must restrict our probability space to a strict subset of $\powerset{\samplespace}$.
More precisely, the sets $\sB \subseteq \samplespace$ for which a probability (or more generally, a \emph{volume}) can not be defined are called \emph{non-measurable sets}. 
An example of such sets is Vitali's sets.\footnote{See \url{https://en.wikipedia.org/wiki/Non-measurable_set} and \url{https://en.wikipedia.org/wiki/Vitali_set}.}
See also Appendix A.2 in \citet{Durrett2019}.

Later, we will be interested in modeling probability spaces over sets of (infinite) sequences.
By virtue of a theorem due to Carath\'{e}odory, there is a natural way to construct such a probability space for sequences (and many other spaces) that behaves in accordance with our intuition, as we will clarify later. 
Here, we shall lay out a few other necessary definitions.
\begin{definition}
$\eventspaceA \subseteq \powerset{\samplespace}$ is called an \defn{algebra} (or field) over $\samplespace$ if
\begin{enumerate}[label=\arabic*)]
    \item $\samplespace\in\eventspaceA$,
    \item if $\sE\in\eventspaceA$, then $\setcomplement{\sE} \in \eventspaceA$,
    \item if $\sE_1,\sE_2 \in \eventspaceA$, then $\sE_1 \cup \sE_2 \in \eventspaceA$.
\end{enumerate}
\end{definition}

\begin{definition}
Let $\eventspaceA$ be an algebra over some set $\samplespace$. A \defn{probability pre-measure} over $(\samplespace,\eventspaceA)$ is a function $\probfunction_0:\eventspaceA \to [0,1]$ such that
\begin{enumerate}[label=\arabic*)]
    \item $\probfunction_0(\samplespace)=1$,
    \item if $\sE_1,\sE_2,\dots$ is a (countable) sequence of disjoint sets in $\eventspaceA$ \emph{whose (countable) union is also in} $\eventspaceA$, then $\probfunction_0(\cup_{\idx=1}^\infty \sE_\idx)=\sum_{\idx=1}^{\infty} \probfunction_0(\sE_\idx)$.
\end{enumerate}
\end{definition}

Note that the only difference between a $\sigma$-algebra (\cref{def:sigma-algebra}) and an algebra is that condition 3 is weakened from countable to finite, and the only difference between a probability measure (\cref{def:probability-measure}) and a pre-measure is that the latter is defined with respect to an algebra instead of a $\sigma$-algebra.

The idea behind Carath\'eodory's extension theorem is that there is often a simple construction of an algebra $\eventspaceA$ over $\samplespace$ such that there is a natural way to define a probability pre-measure.
One can then \emph{extend} this probability pre-measure to a probability measure that is both minimal and unique in a precise sense. 
For example, the standard Lebesgue measure over the the real line can be constructed this way.

Finally, we define random variables.
\begin{definition}[Random variable]\label{def:rv}
A mapping $\rx:\samplespace\to S$ between two measurable spaces $(\samplespace, \eventspace)$ and $(S,\mathcal{S})$ is an $(S,\mathcal{S})$-valued \defn{random variable}, or a measurable mapping, if, for all $\sB \in \mathcal{S}$,
\begin{equation}
    \inv{\rx}(\sB)\defeq \{\omega\in\samplespace:\rx(\omega)\in \sB\} \in \samplespace.
\end{equation}
\end{definition}

Any measurable function (random variable) induces a new probability measure on the \emph{output} $\sigma$-algebra based on the one defined on the original $\sigma$-algebra.
This is called the \defn{pushforward measure} (cf. \S2.4 in \citealp{Tao2011}), which we will denote by $\pushfwdMeasure$, given by
\begin{align} 
    \pushfwdMeasure\left(\rx\in \sE\right)\defeq\probfunction\left(\inv{\rx}\left(\sE\right)\right),
\end{align}
that is, the probability of the result of $\rx$ being in some event $\sE$ is determined by the probability of the event of all the elements which $\rx$ maps into $\sE$, i.e., the pre-image of $\sE$ given by $\rx$.

\ryan{Insert some random variable examples.}    